% Generated draft: requires mathematical and source review.
\documentclass[11pt]{article}
\usepackage[margin=25mm]{geometry}
\usepackage{fontspec}
\setmainfont{Latin Modern Roman}
\usepackage{amsmath,amssymb,mathtools,unicode-math}
\setmathfont{Latin Modern Math}
\usepackage{enumitem,needspace,xurl,hyperref}
\setlength{\parindent}{0pt}
\setlength{\parskip}{5pt}
\setlength{\emergencystretch}{3em}
\newcommand{\N}{\mathbb{N}}
\newcommand{\Z}{\mathbb{Z}}
\newcommand{\Q}{\mathbb{Q}}
\newcommand{\R}{\mathbb{R}}
\newcommand{\C}{\mathbb{C}}
\newcommand{\F}{\mathbb{F}}
\newcommand{\A}{\mathbb{A}}
\newcommand{\PP}{\mathbb{P}}
\newcommand{\E}{\mathbb{E}}
\newcommand{\eps}{\varepsilon}
\newcommand{\dd}{\,\mathrm{d}}
\begin{document}
\section*{UnsolvedMath ID: 10400198}
\textbf{Source problem code:} catalogue entry
\subsection{Definitions and mathematical statement}
Let \(\mathcal H\) denote the integral homology \(3\)-spheres in the convention intended by the source: closed connected \(3\)-manifolds \(M\) with \(H_*(M;\Z)\cong H_*(S^3;\Z)\), taken up to the appropriate equivalence. In the usual oriented smooth formulation, that equivalence is orientation-preserving diffeomorphism. Let
\[
 Z_{\mathrm{LMO}}(M)\in\widehat{\mathcal A}(\varnothing;\Q)
 =\prod_{d\geq0}\mathcal A(\varnothing;\Q)^{(d)}
\]
be the full LMO invariant, with a fixed consistent normalization. Here \(\mathcal A(\varnothing;\Q)\) is the graded rational space of closed vertex-oriented trivalent Jacobi diagrams modulo antisymmetry and IHX, and the hat denotes degree completion.

\emph{Conjecture.} The LMO invariant distinguishes integral homology \(3\)-spheres. In the conventional oriented formulation, for every pair \(M,N\in\mathcal H\),
\[
 Z_{\mathrm{LMO}}(M)=Z_{\mathrm{LMO}}(N)
 \quad\Longrightarrow\quad
 M\cong N\text{ by an orientation-preserving diffeomorphism}.
\]
Equality means equality of the complete series in every degree, not of a single finite truncation. The source cross-refers to Conjecture 10.9; its statement is absent. The excerpt does not explicitly specify orientation, equivalence, or normalization conventions, so these must be checked against the source before adopting the conventional formulation as definitive.
\subsection{Short English statement}
The conjecture says that the complete LMO series is injective on integral homology \(3\)-spheres: two spheres with identical invariants in every degree should be equivalent. It does not assert a fixed finite degree sufficient to distinguish all pairs.
\subsection{Sources}
\begin{itemize}
\item Local UnsolvedMath registry, permanent ID 10400198, source code catalogue entry, ``Conjecture 11.2 --- The LMO invariant distinguishes integral homology 3-spheres.'' No individual proposer is supplied.
\item Ohtsuki, \emph{Problems on invariants of knots and 3-manifolds} (2002), Conjecture 11.2, PDF page 155, with cross-reference to Conjecture 10.9. Supplied URL: \url{https://msp.org/gtm/2002/04/gtm-2002-04-024s.pdf}; recorded access 2026-07-29; extraction pdf-font-structure. The collection distinguishes 231 numbered bold problem, conjecture, and question headings.
\item The registry supplies status open and research\_classification OPEN-TRIAGE; its summary begins ``Classification: OPEN,'' calls the problem major and open, and describes it as equivalent to Conjecture 10.9. Those are supplied assessments, not an independent present-day audit. Current primary-literature status and redistribution terms for the publicly accessible author/publisher copy remain NEEDS\_REVIEW.
\end{itemize}
\end{document}
